\section{The large-degree range and finite reductions}\label{sec:reduction}

The proof has a representation-theoretic reduction and a finite
multiplicity comparison. We first specify the canonical map supplied by
the companion and reduce the remaining comparison to partitions with at
most six parts. A chart of the same invariant ring gives one of the finite
certificates. The final subsection proves the independent large-degree
bound used with the extended computation.

\subsection{The multiplication map}
For the moment let \(a\ge2\), and set
\[
 P=(\Sym V)^{\otimes a},\qquad
 W_j=P_{(j,\ldots,j)},\qquad
 R_j=W_j^{\SSS_a},\qquad R=\bigoplus_{j\ge0}R_j.
\]
The symmetric group permutes the tensor factors. We identify
\(\Sym^a W\) with \((W^{\otimes a})^{\SSS_a}\) by the normalized average
\[
 w_1\cdots w_a\longmapsto
 \frac1{a!}\sum_{\sigma\in\SSS_a}
 w_{\sigma(1)}\otimes\cdots\otimes w_{\sigma(a)}.
\]
Its inverse is the restriction of the symmetric-product quotient.
Let \(A\subset R\)
be the graded subalgebra generated by \(R_1=\Sym^a V\).
For \(v\in V\), put \(z(v)=v^{\otimes a}\in R_1\).
Multiplication defines the canonical Foulkes--Howe map
\begin{equation}\label{eq:multmap}
       \mu_{a,b,V}:\Sym^b(\Sym^a V)\longrightarrow
       R_b=\Sym^a(\Sym^b V),
\end{equation}
whose image is \(A_b\). Thus \(\mu_{a,b,V}=\Psi_{b,a,V}\) in the
notation of the introduction. It is equivariant, so \(A_b=R_b\) suffices
for the desired embedding, by complete
reducibility. This is the equal-multidegree invariant construction
underlying the Chow normalization; see \cite[Section~2]{RaicuSamWeyman2022}.
We will use its explicit polynomial realization.

\begin{corollary}[The main large-degree range]\label{cor:quadratic-range}
For every finite-dimensional complex \(V\) and every \(b\ge30\), there
is a \(\GL(V)\)-equivariant injection
\(\Sym^6(\Sym^b V)\hookrightarrow\Sym^b(\Sym^6 V)\).
\end{corollary}
\begin{proof}
The map in~\eqref{eq:multmap}, with exactly the averaging convention
above, is the canonical multiplication map of
\cite[Theorem~1]{OpenAIQuadratic2026}. That theorem gives \(A_b=R_b\)
for \(a=6\) and \(b\ge6(6-1)=30\), including dimension zero.
In positive dimension its kernel has an invariant complement, by
complete reducibility of complex polynomial \(\GL(V)\)-representations.
Restriction to this complement is an isomorphism onto \(R_b\);
its inverse gives an equivariant section and hence the stated injection.
There is no assertion that this section is a prescribed canonical map.
\end{proof}

\subsection{Six variables and rectangular shifts}
For the finite comparison take \(a=6\). If \(b\ge0\) and
\(\lambda\) is a partition of \(6b\), the Cauchy decomposition of
\(P=\Sym(V\otimes\CC^6)\)
\cite[Exercise~6.11(b)]{FultonHarris1991} gives the source multiplicity
\begin{equation}\label{eq:label-multiplicity}
 F(\lambda)=\dim\left((\Schur_\lambda\CC^6)_{(b,\ldots,b)}^{\SSS_6}\right).
\end{equation}
In particular it vanishes for partitions of length greater than six.
Stable Schur coefficients allow us to work in \(V=\CC^6\): extra
variables are specialized to zero, while longer target constituents
already have nonnegative multiplicities. This covers every dimension.

Put \(S=\Sym(\Sym^6\CC^6)\) with its ordinary symmetric-algebra grading,
and let \(Q(\lambda)\) be the multiplicity of \(\Schur_\lambda\CC^6\)
in \(S_b\). Pad every partition by zeros to six parts and write
\[
 d_i=\lambda_i-\lambda_{i+1},\quad \lambda_7=0,\qquad
 \lambda(d)_i=\sum_{j=i}^6d_j,\qquad
 w(d)=\sum_{i=1}^6 i\,d_i.
\]
Also put \(t_s(d)=\sum_{j>s}\lambda(d)_j\). We abbreviate
\(F(d)=F(\lambda(d))\) and \(Q(d)=Q(\lambda(d))\) when \(6\mid w(d)\);
then \(b=w(d)/6\). Gap-vector inequalities are componentwise, and \(e_i\)
denotes the \(i\)-th unit vector in gap coordinates.

\begin{lemma}\label{lem:rectangles}
Let \(c,d\in\NN^6\) satisfy \(6\mid w(c)\) and \(6\mid w(d)\).
If \(Q(c)>0\), then \(Q(d+c)\ge Q(d)\). In particular
\[
          Q(d+6e_i)\ge Q(d)\quad(1\le i\le6),
          \qquad F(d+6e_6)=F(d).
\]
\end{lemma}
\begin{proof}
Multiplication by one nonzero highest-weight polynomial of weight
\(\lambda(c)\) injects the highest-weight space at \(\lambda(d)\)
into that at their sum, because \(S\) is a domain.
This is the usual highest-weight multiplication argument; compare
\cite[preprint, Lemma~2.2]{BurgisserIkenmeyerPanova2019}.

The rectangular occurrence \(Q(6e_i)>0\) is a special case of the
even-partition theorem of B\"urgisser, Christandl, and Ikenmeyer
\cite[Section~1.1, Theorem]{BurgisserChristandlIkenmeyer2011}.
We give the determinant-tensor construction of
\cite[preprint, Corollary~4.8]{BurgisserIkenmeyerPanova2019}
in the form needed here. Take six copies of the determinant tensor on
the first \(i\) basis vectors. Regroup the tensor positions into \(i\)
blocks, each taking the same position from every copy, and symmetrize
to \(\Sym^i(\Sym^6\CC^6)\). This is a highest vector of weight
\((6,\ldots,6)\) with \(i\) parts. Pair each block with the same sum of
pure sixth powers of the first \(i\) dual basis vectors. A surviving
term uses the same permutation in all six determinant copies, so its
sign is \(+1\); such terms exist. The vector is nonzero.

Adding \(6e_6\) adds six to every partition part.
In~\eqref{eq:label-multiplicity}, this tensors with \(\det^6\),
shifts the equal label weight by six, and is trivial on permutation
matrices. Thus \(F(d+6e_6)=F(d)\).
\end{proof}

It follows that, after checking all partitions for \(6\le b\le25\),
induction in \(26\le b\le29\) need only treat \(d_6<6\).
If \(d_6\ge6\), subtract \(6e_6\); the degree becomes \(b-6\ge20\),
and Lemma~\ref{lem:rectangles} imports the earlier comparison.
The same induction works throughout \(26\le b\le149\), the full
range retained in the programs for the independent route.

\subsection{A chart certificate for the first row}
We return to arbitrary \(a\ge2\) and \(V\) for the chart calculation.
Its factor bound will show, at \(a=6\), that the multiplication image
contains every source highest-weight space with \(\lambda_1\ge5b\).
It will also provide the independent large-degree bound in the next
subsection. Dimension zero is immediate, so assume \(V\ne0\).

Fix \(0\ne v\in V\), use it as the first basis vector, and write
\(x_{i0},x_{i1},\ldots\) for the corresponding polynomial variables in
row \(i\) of \(P\). Put \(z=z(v)=\prod_i x_{i0}\).
After inverting \(z\), every \(x_{i0}\) is invertible. Write
\(Y_i=(x_{i1}/x_{i0},x_{i2}/x_{i0},\ldots)\).

\begin{lemma}\label{lem:chart}
For every integer \(t\ge0\), every row-permutation-invariant polynomial
in the lists \(Y_i\) of total degree at most \(t\) is a linear combination of products of at
most \(t\) elements of \(z^{-1}R_1\).
Consequently, for integers \(k,L\ge0\) and \(f\in R_k\),
\[
             z^Lf\in A_{k+L}\qquad\text{if }k+L\ge ak.
\]
\end{lemma}
\begin{proof}
The multisymmetric generators are described in
\cite[Remark~2.9]{RaicuSamWeyman2022}; we include the factor-length
argument needed here.
For the general multisymmetric generator theorems, see
\cite[Theorem~1]{Vaccarino2005} and \cite[Corollary~8.4]{Rydh2007}.
For each linear form \(\ell\) in one list, the coefficients of
\[
 \frac{z(v+tu)}{z}=\prod_{i=1}^a(1+t\ell(Y_i))
\]
show that \(e_h(\ell(Y_1),\ldots,\ell(Y_a))\), and all their
polarizations, lie in \(z^{-1}R_1\), for \(1\le h\le a\).
Here \(u\) ranges over the span of the remaining basis vectors.

Give \(e_h\) weight \(h\). Newton identities express
\(\sum_i\ell(Y_i)^d\) as a polynomial in \(e_1,\ldots,e_a\)
of weighted degree \(d\), including when \(d>a\). Each product has
at most \(d\) factors. Polarizing in \(\ell\) therefore expresses
\[
                         p_m=\sum_i m(Y_i)
\]
for every degree-\(d\) monomial \(m\) using at most \(d\) polarized
elementary factors.

Invariant monomial orbit sums are scalar multiples of expressions
\[
        \sum_{\substack{i_1,\ldots,i_h\\\text{distinct}}}
                    \prod_{\nu=1}^h m_\nu(Y_{i_\nu}),
\]
where the \(m_\nu\) have positive degree. Inclusion-exclusion over
collisions of the indices expresses this in products of \(p_m\)'s;
a collision replaces monomials by their product and preserves total
degree. Thus total degree at most \(t\) requires at most \(t\) elementary
factors. Constants require no factors.

For \(f\in R_k\), the polynomial \(f/z^k\) is invariant and has total
list degree at most \(ak\). A term with \(m\le ak\) ratio factors has
the form \(z^{-m}g_1\cdots g_m\), with \(g_\nu\in R_1\).
Multiplication by \(z^{k+L}\) gives an element of \(A_{k+L}\) when
\(k+L\ge ak\). Equality in the localization descends to \(P\), a domain.
\end{proof}

\begin{corollary}[First-row certificate]\label{cor:first-row}
For every integer \(b\ge1\) and partition \(\lambda\) of \(6b\) with
\(\sum_{i\ge2}\lambda_i\le b\), one has
\[
             [s_\lambda]h_6[h_b]\le [s_\lambda]h_b[h_6].
\]
\end{corollary}
\begin{proof}
Partitions of length greater than six have zero source multiplicity, so
work in \(V=\CC^6\) with its standard ordered basis and put \(a=6\).
Choose \(v\) to be the first basis vector in the chart above.
For every weight-\(\lambda\) polynomial \(f\in R_b\), the total
list degree of \(f/z^b\) is
\(t=\lambda_2+\cdots+\lambda_6\).
The first assertion of Lemma~\ref{lem:chart} writes this ratio as a
linear combination of terms \(z^{-m}g_1\cdots g_m\), with
\(m\le t\le b\) and \(g_\nu\in R_1\). Thus every corresponding
term of \(f\) is
\[
                     z^{b-m}g_1\cdots g_m\in A_b.
\]
The entire weight-\(\lambda\) space of \(R_b\) therefore lies in
\(A_b\). Since \(A_b\) is a \(\GL(V)\)-submodule of \(R_b\), their
highest-weight spaces at \(\lambda\) are equal: intersect the common
weight space with the kernels of the positive simple-root operators.
The surjection onto \(A_b\) induced by~\eqref{eq:multmap} then gives
\(Q(\lambda)\ge F(\lambda)\).
\end{proof}

\subsection{An independent large-degree bound}
We now prove surjectivity of \(\mu_{a,b,V}\) for
\(b\ge a(a-1)^2\), without using the companion. The chart will clear
denominators uniformly once we know that \(R\), as an \(A\)-module,
is generated in degrees below \(a\). Polarization supplies that
finite-generation statement and then turns the chart calculation into
surjectivity in large degree. We continue with arbitrary \(a\ge2\);
the zero-dimensional case is immediate.

\begin{lemma}[Polarization]\label{lem:polarization}
For integers \(L\ge1\) and \(j\ge aL\),
\[
       W_j=\sum_{v\in V}z(v)^L W_{j-L},
       \qquad
       R_j=\sum_{v\in V}z(v)^L R_{j-L}.
\]
The sums mean linear spans.
\end{lemma}
\begin{proof}
The operator argument is related to the successive weight shifts in
\cite[Section~5]{Ikenmeyer2015}. Polarize the degree-\(aL\) polynomial
\(v\mapsto z(v)^L\), and multiply, to obtain a linear map
\[
                 \Sym^{aL}V\otimes W_{j-L}\longrightarrow W_j.
\]
Its image is the first displayed span. Identify \(W_j^*\) with
polynomials \(f(x_1,\ldots,x_a)\), separately homogeneous of degree \(j\).
Introduce one auxiliary vector variable \(y\). The dual map is, up to a nonzero scalar,
\[
                f\longmapsto\prod_{i=1}^a
                       (y\cdot\partial_{x_i})^L f.
\]
Indeed, evaluation of the dual pairing extracts the coefficient of
\(t_1^L\cdots t_a^L\) in \(f(x_1+t_1y,\ldots,x_a+t_ay)\).

On the polynomial space of total degree \(aj\) in
\(x_1,\ldots,x_a,y\), the operators
\[
 D_i=y\cdot\partial_{x_i},\quad
 E_i=x_i\cdot\partial_y,\quad
 H_i=\deg(x_i)-\deg(y)
\]
form an \(\mathfrak{sl}_2\), with \(D_i\) lowering weight by two.
The other lists are spectator variables. The standard weight strings
are described in \cite[Section~11.1, (11.5)]{FultonHarris1991}.
In an irreducible module of highest weight \(n\), a weight \(w\ge L\)
is at distance \((n+w)/2\ge L\) from the bottom. Therefore \(D_i^L\)
is injective on the full weight-\(w\) space of any finite-dimensional
module. After the first \(i-1\) shifts, the relevant weight is
\(j-(i-1)L\ge L\). Every shift is injective on a space containing the
preceding image, so their product is injective. Duality proves the
first identity. Average its coefficients over \(\SSS_a\); each
\(z(v)^L\) is invariant, proving the second identity.
\end{proof}

\begin{corollary}\label{cor:module-generation}
The \(A\)-module \(R\) is generated by \(R_0,\ldots,R_{a-1}\).
\end{corollary}
\begin{proof}
Lemma~\ref{lem:polarization} with \(L=1\) gives
\(R_j=R_1R_{j-1}\) for \(j\ge a\); apply induction on \(j\).
\end{proof}

\begin{proposition}[Independent chart bound]\label{prop:large}
For \(a\ge2\), one has \(R_b=A_b\) whenever
\(b\ge a(a-1)^2\).
\end{proposition}
\begin{proof}
Take \(L=(a-1)^2\), so that \(L\ge(a-1)k\) for every
\(0\le k<a\). Lemma~\ref{lem:chart} therefore applies to each of
these degrees. For any fixed nonzero \(v\), multiply an \(A\)-linear
expression in the generators of Corollary~\ref{cor:module-generation}
by \(z(v)^L\); this gives \(z(v)^LR\subset A\).
The same assertion is trivial for \(v=0\). If \(b\ge aL\),
Lemma~\ref{lem:polarization} now yields
\[
                  R_b=\sum_v z(v)^L R_{b-L}\subseteq A_b.
\]
The opposite inclusion is definitional.
\end{proof}

For \(a=6\), Proposition~\ref{prop:large} supplies the independent
large-degree range \(b\ge150\). Combined with the retained finite
verification through \(b=149\), it proves Theorem~\ref{thm:main}
without the quadratic-stabilization input. The main proof needs only
\(6\le b\le29\), since Corollary~\ref{cor:quadratic-range} applies
from \(b=30\).
